% Propietario conservado: /Users/ruben/Documents/New project/output/EXTREMA_MEDIA_RAZON_RECIPROCIDAD_ES_EN_20260918/ARTICULO/ES/source/nuclear/06.tex
\chapter{Affine holonomy and memory symmetries}
\label{x_reciprocidad:nuclear:6}
\section{Effective composition of memory}

For an admissible emission \(\epsilon\), the construction in \cref{x_reciprocidad:iv:cont:concatenacion} defines

\begin{equation}
\mathsf U_\epsilon(m)=C_\epsilon m+\kappa_\epsilon,
\qquad
\kappa_{\epsilon_2\epsilon_1}
=\kappa_{\epsilon_2}+C_{\epsilon_2}\kappa_{\epsilon_1}.
\label{x_reciprocidad:nuc06:eq:1}
\end{equation}

The domain is the memory module at the source; the codomain is that of the next emission. The product preserves order. On a path \(\gamma\) of length \(r\), one obtains

\begin{equation}
C_\gamma=C_r\cdots C_1,\qquad
\kappa_\gamma=\sum_{j=1}^r C_r\cdots C_{j+1}\kappa_j,
\qquad
m_\gamma=C_\gamma m_0+\kappa_\gamma.
\label{x_reciprocidad:nuc06:eq:2}
\end{equation}

The empty product in the sum is the identity. The proof is induction using \eqref{x_reciprocidad:nuc06:eq:1}: appending the next emission multiplies the preceding record by its linear transport and adds its increment. The positional boundary may cancel while the transported sum in \eqref{x_reciprocidad:nuc06:eq:2} remains.

\begin{proposition}[Memory class of a return]

Let \(M\) be an abelian module, and consider a path closed in the base whose memory transport is \(H(m)=Cm+\kappa\), with \(C\) an automorphism of \(M\). Its return class is

\begin{equation}
\mathfrak m(H)=[\kappa]\in M/(I-C)M.
\label{x_reciprocidad:nuc06:eq:3}
\end{equation}

A memory-coordinate change \(m'=Rm+b\), with \(R\) invertible, transforms the return into

\begin{equation}
C'=RCR^{-1},\qquad
\kappa'=R\kappa+(I-C')b.
\label{x_reciprocidad:nuc06:eq:4}
\end{equation}

The map induced by \(R\) takes \([\kappa]\) to \([\kappa']\). The transport \(H\) has a fixed point exactly when its class \eqref{x_reciprocidad:nuc06:eq:3} is zero.
\end{proposition}

\begin{proof}
Substituting \(m=R^{-1}(m'-b)\) into \(H\) gives \eqref{x_reciprocidad:nuc06:eq:4}. Since \(R(I-C)=(I-C')R\), the map \(R\) induces an isomorphism between the quotients. The term \((I-C')b\) represents zero in the target quotient. Finally, the fixed-point equation \(Cm+\kappa=m\) is equivalent to \((I-C)m=\kappa\). This proves the final assertion.
\end{proof}

Formula \eqref{x_reciprocidad:nuc06:eq:3} determines whether a change of origin removes the increment or merely redistributes it among coordinates. In particular, when \(C=I\), the class is the vector \(\kappa\) itself. A nonzero translation remains nonzero under every invertible affine change.

\begin{proposition}[Return symmetry versus a change of representation]

An affine transformation \(A(m)=Rm+b\) commutes with \(H\) exactly when

\begin{equation}
RC=CR,\qquad (I-C)b=(I-R)\kappa.
\label{x_reciprocidad:nuc06:eq:5}
\end{equation}
\end{proposition}

\begin{proof}
The linear parts of \(AH\) and \(HA\) are \(RC\) and \(CR\); their constant parts are \(R\kappa+b\) and \(Cb+\kappa\). Equating them gives \eqref{x_reciprocidad:nuc06:eq:5}. If \(H\) is merely transported through \eqref{x_reciprocidad:nuc06:eq:4}, commutation is not required: the same law is obtained in another representation. By contrast, a symmetry of the fixed return requires both equalities in \eqref{x_reciprocidad:nuc06:eq:5}.
\end{proof}

This precisely distinguishes covariance of the law from the stabilizer of a specific memory. The distinction is relevant to studying variations in state symmetry without modifying the composition laws.

\section{Application to the effective record of the observable revolution}

Section~\ref{x_reciprocidad:vii:sec:retorno-memoria-gravedad} defines two ordered counters: \(g\) counts orientation reversals, and \(s\) counts effective sheet changes. Their canonical count per 27-step block is

\[
c_{27}=(1,18).
\]

The complete observable phase returns after four such blocks. In the path groupoid, the record module is \(M=\mathbb Z^2\), and

\begin{equation}
c_{108}=(4,72)=4v,\qquad v=(1,18),\qquad
H(g,s)=(g+4,s+72).
\label{x_reciprocidad:nuc06:eq:6}
\end{equation}

Identity \eqref{x_reciprocidad:nuc06:eq:6} is a corpus result used here; the verifier checks its consequences rather than recounting the original emissions. The symbol \(c_{108}\) denotes this two-counter increment. It differs from the dodecaphasic seal \(K\) and from the central cocycle of the Heisenberg representation.

\begin{theorem}[Complete record invariants under full revolutions]

The orbits of \(H\) and its inverse on \(\mathbb Z^2\) are classified exactly by

\begin{equation}
\boxed{\quad \mathcal I(g,s)=s-18g,\qquad \nu(g,s)=[g]_4.\quad}
\label{x_reciprocidad:nuc06:eq:7}
\end{equation}

Consequently,

\begin{equation}
\mathbb Z^2/\mathbb Z(4,72)\cong\mathbb Z\oplus\mathbb Z/4\mathbb Z.
\label{x_reciprocidad:nuc06:eq:8}
\end{equation}
\end{theorem}

\begin{proof}
Adding \((4,72)\) preserves both expressions in \eqref{x_reciprocidad:nuc06:eq:7}. If two records have the same invariants, then \(g'-g=4N\) for an integer \(N\), and \(s'-s=18(g'-g)=72N\); thus they differ by exactly \(N(4,72)\). The map in \eqref{x_reciprocidad:nuc06:eq:7} is surjective: given an integer \(j\) and \(r\) modulo four, choose a representative \(g\) of \(r\) and set \(s=j+18g\). This proves \eqref{x_reciprocidad:nuc06:eq:8}.
\end{proof}

The description covers every orbit of this memory reader. Orbits of the complete enriched state additionally preserve their remaining coordinates. For a forward-directed prolongation, a later record is obtained when the preceding integer N is nonnegative; the bilateral equivalence in \eqref{x_reciprocidad:nuc06:eq:8} belongs to the groupoid.

The quantity \(\mathcal I\) is the unique primitive integer linear charge, up to sign, invariant under these translations. Indeed, for \(L(g,s)=ag+bs\), invariance requires \(4a+72b=0\), so \((a,b)=b(-18,1)\). A revolution accumulates longitudinal memory and preserves a transverse relation determined by the increment.

The factor four has a precise origin here: it is \(\gcd(4,72)\) and, in the preceding count, the number of 27-step blocks completing the observable return. The class \(\nu\) preserves the information lost when the increment is divided by four and only the primitive direction \(v\) is retained.

\begin{corollary}[Joint block and phase record]

Write the phase of these four blocks as \(a\in\mathbb Z/4\mathbb Z\). The update

\[
(a,g,s)\longmapsto(a+1,g+1,s+18)
\]

preserves \((s-18g,[g-a]_4)\). Both data classify its bilateral orbits: the difference \(N=g'-g\) simultaneously satisfies \(a'-a=[N]_4\) and \(s'-s=18N\). This brings together memory advance and phase return without replacing either by the other.
\end{corollary}

\section{Parabolic generator determined by memory}

The ordered pair of counters equips the module \(M\) with the unimodular alternating form

\[
\Omega_M((g,s),(g',s'))=gs'-sg'.
\]

Its domain is the record module. This form is not identified with a dimensional physical area or with the form on the APP torus, whose domain differs. The choice of order \((g,s)\) fixes its orientation.

\begin{theorem}[Unimodular stabilizer of the increment]

Define, directly from the primitive vector \(v\) in \eqref{x_reciprocidad:nuc06:eq:6},

\begin{equation}
N_v m=v\,\Omega_M(v,m)=v\mathcal I(m),\qquad
N_v=\begin{pmatrix}-18&1\\-324&18\end{pmatrix}.
\label{x_reciprocidad:nuc06:eq:9}
\end{equation}

Then \(N_v^2=0\), \(N_v\ne0\), and

\begin{equation}
\boxed{\quad
\operatorname{Stab}_{SL_2(\mathbb Z)}(c_{108})
=\{I+nN_v:n\in\mathbb Z\}.
\quad}
\label{x_reciprocidad:nuc06:eq:10}
\end{equation}

Every transformation in \eqref{x_reciprocidad:nuc06:eq:10} preserves \(\Omega_M\), the increment \(c_{108}\), and the charge \(\mathcal I\). For nonzero \(n\), it is a nonidentity unipotent matrix of trace two, of parabolic type.
\end{theorem}

\begin{proof}
The identity \(\Omega_M(v,v)=0\) implies \(N_v^2m=v\Omega_M(v,v)\Omega_M(v,m)=0\). Its upper-right entry is one, so it is nonzero. The matrix

\[
B=\begin{pmatrix}1&0\\18&1\end{pmatrix}
\]

has determinant one, maps the first basis vector to \(v\), and satisfies

\begin{equation}
N_v=BNB^{-1},\qquad N=\begin{pmatrix}0&1\\0&0\end{pmatrix}.
\label{x_reciprocidad:nuc06:eq:11}
\end{equation}

An integer automorphism fixes \(4v\) exactly when it fixes \(v\), because \(M\) is torsion-free. A determinant-one matrix fixing the first basis vector necessarily has the form \(\begin{pmatrix}1&n\\0&1\end{pmatrix}\). Conjugation by \(B\) proves \eqref{x_reciprocidad:nuc06:eq:10}. The rest follows from this expression or directly from \eqref{x_reciprocidad:nuc06:eq:9}.
\end{proof}

Formula \eqref{x_reciprocidad:nuc06:eq:9} makes the construction independent of the second vector chosen to complete \(v\) to an oriented unimodular basis: any other complement differs by an integer multiple of \(v\) and produces the same \(N_v\). Under an oriented transformation \(R\in SL_2(\mathbb Z)\), one has \(N_{Rv}=RN_vR^{-1}\).

The action can be written particularly clearly as

\begin{equation}
(I+nN_v)m=m+n\mathcal I(m)v.
\label{x_reciprocidad:nuc06:eq:12}
\end{equation}

The transverse charge determines how far this symmetry shifts along the memory direction. Summing revolutions produces the displacement \(4Nv\); stabilizing symmetries preserve this advance law.

On the quotient \eqref{x_reciprocidad:nuc06:eq:8}, the transformation \(I+nN_v\) acts as \((\mathcal I,[g]_4)\mapsto(\mathcal I,[g+n\mathcal I]_4)\). It therefore preserves the family of orbits and may permute its residue classes; \(\nu\) is an invariant of the return \(H\), not of every stabilizing transformation. This difference follows directly from \eqref{x_reciprocidad:nuc06:eq:12}.

Conjugation \eqref{x_reciprocidad:nuc06:eq:11} supplies an explicit algebraic map from the parabolic TRIT model, defined in \cref{x_reciprocidad:euler-trit-generadores}, to the record module extended to real scalars. One also obtains \(\exp(tN_v)=I+tN_v\) and \(B\exp(tN)=\exp(tN_v)B\). This intertwining concerns operators: the particular motion selected by TPK also uses its emissions and calendar. The complete arithmetic group in \eqref{x_reciprocidad:nuc06:eq:10} fixes this reader; the admissible symmetry group of the total state is obtained by intersecting it with the actions preserving the remaining data and relations.

In the translated system \eqref{x_reciprocidad:nuc06:eq:6}, affine changes \(m\mapsto Rm+b\), with \(R\) in \eqref{x_reciprocidad:nuc06:eq:10} and arbitrary \(b\), commute with \(H\). Those additionally preserving the fixed initial record must satisfy \(Rm_0+b=m_0\). Distinguishing the two conditions separates symmetry of the law from symmetry of a marked trajectory.

\begin{theorem}[Variational realization of the charge and stabilizer]
\label{x_reciprocidad:nuc06:variacional}

The scalar extension of the record, \(M_{\mathbb R}=M\otimes_{\mathbb Z}\mathbb R\), preserves the generated coordinates \((g,s)\) and the form \(\Omega_M=dg\wedge ds\). It is a subsequent realization of the integer module. With the convention \(\iota_{X_H}\Omega_M=dH\), one has

\begin{equation}
X_{\mathcal I}=v,\qquad
X_{\frac12\mathcal I^2}=\mathcal I v=N_vm.
\label{x_reciprocidad:nuc06:eq:12a}
\end{equation}

Thus, the constructed translation and parabolic family admit explicit Hamiltonian generators. The flow of \(\mathcal I\) is \(m\mapsto m+uv\); at \(u\) equal to four it reproduces \eqref{x_reciprocidad:nuc06:eq:6}. The flow of \(\mathcal I^2/2\) is \(m\mapsto(I+uN_v)m\). The parameter \(u\) corresponds to this continuous realization; the discrete return step and intermediate emissions retain their own TPK calendar.
\end{theorem}

\begin{proof}
For a vector \(X=(X_g,X_s)\), \(\iota_X(dg\wedge ds)=X_g\,ds-X_s\,dg\). Since \(d\mathcal I=ds-18\,dg\), the first field is \((1,18)=v\). The chain rule gives the second. The quantity \(\mathcal I\) remains constant along both fields because \(d\mathcal I(v)=0\). The second flow therefore integrates \(dm/du=\mathcal I(m_0)v\), equivalent to \(\exp(uN_v)=I+uN_v\). Nondegeneracy of \(\Omega_M\) implies that the Hamiltonian generating \(N_v\) is \(\mathcal I^2/2\) plus an additive constant: its differential is fixed by \(\iota_{N_vm}\Omega_M\).
\end{proof}

The mathematical action

\begin{equation}
\mathscr A[m]=\int_{u_0}^{u_1}
\left[s\frac{dg}{du}-\frac12(s-18g)^2\right]du
\label{x_reciprocidad:nuc06:eq:12b}
\end{equation}

yields, by variation, the equations \(dg/du=\mathcal I\), \(ds/du=18\mathcal I\). The translation \(g\mapsto g+\varepsilon\), \(s\mapsto s+18\varepsilon\) preserves the Hamiltonian and changes the Lagrangian by the total derivative \(d(18\varepsilon g)/du\), for constant \(\varepsilon\). Its Noether charge is

\begin{equation}
\underbrace{s\,\delta g}_{s}
-\underbrace{18g}_{\text{boundary term}}
=s-18g=\mathcal I.
\label{x_reciprocidad:nuc06:eq:12c}
\end{equation}

The variation is interpreted per unit of \(\varepsilon\). The equations directly verify \(d\mathcal I/du=18\mathcal I-18\mathcal I=0\). The discrete translation \eqref{x_reciprocidad:nuc06:eq:6} itself is also realized by the Hamiltonian \(H_{\rm ret}=4\mathcal I\), whose flow at parameter one is \(H\); the two Hamiltonians Poisson-commute. The stabilizer flow \(I+uN_v\) and the return \(m+4v\) are distinct commuting actions: on the line \(\mathcal I=0\), the former fixes every point, while the latter continues to accumulate memory.

More generally, every \(C^1\) Hamiltonian on this plane invariant under translations \(m\mapsto m+uv\) has the form \(F(\mathcal I)\). To prove this, the global coordinates \((g,\mathcal I)\) express the symmetry as translation in \(g\); invariance removes dependence on \(g\) from \(H\). Its field is \(F'(\mathcal I)v\). This classification makes explicit both the common conservation law and the additional choice of specific dynamics.

Action \eqref{x_reciprocidad:nuc06:eq:12b} is dimensionless and belongs to the realized record plane. Its proof supplies an exact variational relation between memory, charge, and stabilizer. Conversion to physical action and dimensional time is performed through the corresponding action and clock readers; these units are not assigned to the counter u.

\section{APP curvature and conservation of the central phase}

Section~\ref{x_reciprocidad:iv:extension-curvatura-mixta} works on \(X_9=(\mathbb Z/9\mathbb Z)^2\). The sheets are \(S(x,y)=x+y\) and \(P(x,y)=xy\). The mixed connections use differences of these potentials, not the potentials themselves as links:

\begin{equation}
F(S,P)=\Delta_x\Delta_yP-\Delta_y\Delta_xS=1,
\qquad F(P,S)=-1\pmod9.
\label{x_reciprocidad:nuc06:eq:13}
\end{equation}

The oriented sum of links around a rectangle is \(ab\) or \(-ab\), depending on order. The alternating area is \(\sigma(u,v)=ad-bc\) for \(u=(a,b)\), \(v=(c,d)\). Section~\ref{x_reciprocidad:iv:weyl-area-extension} represents these transports through \(D(u)=X^aZ^b\), \(ZX=\omega XZ\), with \(\omega\) the ninth root in the subsequent representation. Then

\begin{equation}
D(u)D(v)=\omega^{bc}D(u+v),\qquad
D(u)D(v)D(u)^{-1}D(v)^{-1}=\omega^{-\sigma(u,v)}I.
\label{x_reciprocidad:nuc06:eq:14}
\end{equation}

The position resulting from the commutator returns. The central phase retains oriented area with the sign fixed in \eqref{x_reciprocidad:nuc06:eq:14}. This is retention modulo nine: an area divisible by nine has trivial central phase, although a lift with quotient and record may distinguish its paths. The phase does not replace that lift.

\begin{proposition}[Exact covariance of central transport]

Since two is invertible modulo nine, define

\begin{equation}
W(a,b)=\omega^{ab/2}D(a,b),\qquad
W(u)W(v)=\omega^{-\sigma(u,v)/2}W(u+v).
\label{x_reciprocidad:nuc06:eq:15}
\end{equation}

For \(R\in SL_2(\mathbb Z/9\mathbb Z)\), the map

\begin{equation}
\omega^zW(u)\longmapsto\omega^zW(Ru)
\label{x_reciprocidad:nuc06:eq:16}
\end{equation}

is an automorphism of the central extension. It preserves its products and commutation phases.
\end{proposition}

\begin{proof}
The phase in the product of the two \(W\) is
\(\frac{ab}{2}+\frac{cd}{2}+bc-\frac{(a+c)(b+d)}2=\frac{bc-ad}{2}\). This proves \eqref{x_reciprocidad:nuc06:eq:15}. The identity \(\sigma(Ru,Rv)=\det(R)\sigma(u,v)=\sigma(u,v)\) shows that \eqref{x_reciprocidad:nuc06:eq:16} preserves the law; invertibility of \(R\) supplies the inverse automorphism. If \(\det R=-1\), the corresponding automorphism also changes \(z\mapsto-z\); the area sign and centre are reversed jointly.
\end{proof}

The complete record and its central phase have different types. Covariance preserves the area law while the presentation changes; the affine class in \eqref{x_reciprocidad:nuc06:eq:3} preserves memory content while its origin changes. Both properties act on typed path readers, and their maps must remain separate when composed with the exceptional readout of the complete state.

\begin{proposition}[Persistence under local frames and subdivisions]

Let \(\gamma\) be a loop based at \(x\), with invertible links \(U_\epsilon\). The local change \(U'_\epsilon=G_tU_\epsilon G_s^{-1}\) transforms its holonomy into \(H'_\gamma=G_xH_\gamma G_x^{-1}\): frames at interior vertices cancel successively. If \(H_\gamma=\omega^kI\), the central phase remains exactly the same. For refined links whose product reproduces the preceding link, subdivision likewise preserves \(H_\gamma\), by associativity. On APP plaquettes, cancellation of interior edges realizes this property through finite Stokes.
\end{proposition}

An elementary rectangle has additive Wilson value one and nontrivial central phase. This mixed transport therefore distinguishes actual holonomy from the identity that links of the form \(G_tG_s^{-1}\) would produce. Conservation under subdivision concerns compatible loops and links: adding another area changes the transport.

The integer lift of holonomy distinguishes oriented area from its nonadic remainder: for a rectangle with sides four and seven, \(28=1+9\cdot3\), whereas \(-28=8+9(-4)\) corresponds to the reverse orientation. The residual phase coincides for areas one and twenty-eight, while their quotients distinguish the accumulations. This example exhibits the difference between finite phase and complete record; its area quotient and the temporal counter \(q_{\rm mem}\) retain distinct domains.

\section{Composition between the three regimes and continuity of the balance}

The realizations in \cref{x_reciprocidad:euler-trit-generadores,x_reciprocidad:euler-trit-evaluaciones} construct generators \(J_k^2=-\tau_kI\) on their own fibres, of real dimension two. Take effectively specified links \(B_k:V_k\to V_{k+1}\) preserving the alternating forms, \(B_k^*\Omega_{k+1}B_k=\Omega_k\). The symbol \(*\) in this section means transpose/bilinear pullback. The operators \(E_k=\exp(t_kJ_k)\) preserve \(\Omega_k\): in dimension two, zero trace of the generator implies \(J_k^*\Omega_k+\Omega_kJ_k=0\). The symplectic identity is the hypothesis that would be retained in extending the lemma to higher dimensions.

For the nonadic compression and its recorded complement,

\[
T_k=B_k(8I+E_k)/9,\qquad
D_k=B_k\sqrt8(E_k-I)/9,
\]

algebraic expansion gives

\begin{equation}
\Omega_k=T_k^*\Omega_{k+1}T_k+D_k^*\Omega_{k+1}D_k.
\label{x_reciprocidad:nuc06:eq:17}
\end{equation}

Indeed, the cross terms cancel, and the numerator is \(72\Omega_k+9E_k^*\Omega_kE_k=81\Omega_k\). For \(P_0=I\), \(P_{k+1}=T_kP_k\), successive substitution in path order gives

\begin{equation}
\boxed{\quad
\Omega_0=P_N^*\Omega_NP_N+
\sum_{k=0}^{N-1}(D_kP_k)^*\Omega_{k+1}(D_kP_k).
\quad}
\label{x_reciprocidad:nuc06:eq:18}
\end{equation}

The terminal component and every complement remain present. The identity is oriented and bilinear; positivity results apply in the positive profiles proved in the preceding sections. The conserved form does not require a link between distinct regimes to intertwine their generators: if \(BJ_\tau=J_{\tau'}B\), squaring with invertible \(B\) would impose \(\tau=\tau'\). Conservation of a common structure allows the three regimes to remain distinct.

The complete transport \(U_N=(B_{N-1}E_{N-1})\cdots(B_0E_0)\) maps \(V_0\) to \(V_N\). Matrix monodromy requires a typed closure \(K:V_N\to V_0\); only then is \(KU_N\) an endomorphism whose trace can be computed. Compatible basis changes preserve this trace by conjugation. Phase return, fibre closure, and return of the vector remain different conditions.

If the links have exclusively the form \(B_{yx}=G_yG_x^{-1}\), their product along a path is \(G_{\rm final}G_{\rm inicial}^{-1}\). Closing with the same frame gives the identity. The holonomy in \eqref{x_reciprocidad:nuc06:eq:14} and the increment in \eqref{x_reciprocidad:nuc06:eq:6} exhibit concrete mechanisms that this mere substitution of frames cannot represent by itself.

\section{Compatibility with refinement and arbitrary extent}

The memory action is compatible with a restriction homomorphism of modules \(p:M'\to M\) when

\begin{equation}
pC'=Cp,\qquad p\kappa'=\kappa.
\label{x_reciprocidad:nuc06:eq:19}
\end{equation}

Then \(pH'=Hp\), and \(p\) induces a quotient map

\begin{equation}
M'/(I-C')M'\longrightarrow M/(I-C)M,
\qquad [\kappa']\longmapsto[\kappa].
\label{x_reciprocidad:nuc06:eq:20}
\end{equation}

\begin{proof}
The first equality in \eqref{x_reciprocidad:nuc06:eq:19} maps \((I-C')M'\) into \((I-C)M\); the second fixes the image of the increment. Compatibility of \(H\) follows by substitution. Concatenations and refinements compose the same squares by virtue of law \eqref{x_reciprocidad:nuc06:eq:1}.
\end{proof}

A nonzero return class appearing at a preserved level remains nonzero in every compatible lift: a zero class would project to zero. This statement applies to the system of memories and restrictions satisfying \eqref{x_reciprocidad:nuc06:eq:19}; the joint constructor and the maps of its inverse system are preserved in \cref{x_reciprocidad:iv:continuo-conjunto}.

In the register of \eqref{x_reciprocidad:nuc06:eq:6}, for each integer \(N\) we have \(H^N(m)=m+N(4,72)\). The proof is algebraic for arbitrary \(N\), not an inference from a finite number of iterations. In forward prolongation, each prefix preserves its finite register; the complete history is described by the compatible family of these registers, without replacing its counters with an “infinite” integer quantity.

If complete dynamics has the reader \(q_{\rm mem}\) with \(q_{\rm mem}(\Gamma_9x)=q_{\rm mem}(x)+1\), it likewise cannot have a periodic state of positive period \(N\) on that domain: applying the reader would yield \(q_{\rm mem}(x)=q_{\rm mem}(x)+N\). This corollary uses the complete integer degree; its residual quotients may have finite periods, as appropriate to the observable calendar in the corpus.

\section{Foundational consequence and scope of the composition}

The specific connection obtained is

\begin{equation}
\text{return increment }4v
\ \longrightarrow\ \mathcal I(m)=\Omega_M(v,m)
\ \longrightarrow\ N_vm=v\mathcal I(m)
\ \longrightarrow\ I+nN_v.
\label{x_reciprocidad:nuc06:eq:21}
\end{equation}

The accumulated register determines a conserved charge and a parabolic family stabilizing its direction. Evolution and conservation belong to the same action: \(g\) and \(s\) change, their relation \(s-18g\) remains, and the increment preserves an explicit stabilizer. This makes precise part of the link between memory and symmetry that motivated the investigation.

Equality \eqref{x_reciprocidad:nuc06:eq:18} connects this research with the preceding balances: an evolution may redistribute the form between readout and complements while preserving the total. The APP identities \eqref{x_reciprocidad:nuc06:eq:13}–\eqref{x_reciprocidad:nuc06:eq:16} show, in another concrete reader of the same origin, how transport order produces a central datum even when position returns.

The charge \eqref{x_reciprocidad:nuc06:eq:7} is an algebraic invariant of transport and, through the explicit action \eqref{x_reciprocidad:nuc06:eq:12b}, a Noether charge of the record realization. Theorem~\ref{x_reciprocidad:nuc06:variacional} and its variational development construct this relation, its symmetry, and its boundary term. The preceding physical action and current developments retain their domains and units; transporting them to this realization requires preservation of their dimensional maps.

Composition with dimensional readers must transport the pair \((v,\Omega_M)\) and the charge \(\mathcal I\), preserving their units and the variational current. Every subsequent relation with \(\pi,\phi,e,\alpha\) and \(K\) must arise from this composition; the present formulas specify the information to be preserved.
